\begin{Lema}\label{Lema:Esfera_retracto_de_Rn_menos_origen}
Para cualquier $n\in\mathbb{Z}^+$, $S^{n-1}$ es un retracto de deformación de $\mathbb{R}^n\setminus\{0\}$.
\end{Lema}
\begin{proof}
Supongamos que $\frac{x}{\lVert x\rVert}t+(1-t)x=0$ para algún $(x,t)\in(\mathbb{R}^n\setminus\{0\})\times I$. Entonces
\begin{align*}
\left\lVert\frac{x}{\lVert x\rVert}t+(1-t)x\right\rVert&=0\\
\left\lVert x\left(\frac{t}{\lVert x\rVert}+(1-t)\right)\right\rVert&=0\\
\lVert x\rVert\left|\frac{t}{\lVert x\rVert}+(1-t)\right|&=0\\
\left|t+\lVert x\rVert(1-t)\right|&=0\\
t+\lVert x\rVert(1-t)&=0\\
t+\lVert x\rVert-\lVert x\rVert t&=0\\
t(1-\lVert x\rVert)&=-\lVert x\rVert\\
t\left(\frac{\lVert x\rVert-1}{\lVert x\rVert}\right)&=0.
\end{align*}
Si $t=0$ entonces, como $\frac{x}{\lVert x\rVert}t+(1-t)x=0$, se tiene $x=0$ (contradicción). Si $\frac{\lVert x\rVert-1}{\lVert x\rVert}=0$ entonces $\lVert x\rVert=1$, y sustituyendo en $\frac{x}{\lVert x\rVert}t+(1-t)x=0$ se tiene $x=0$ (contradicción). Por tanto $\frac{x}{\lVert x\rVert}t+(1-t)x\neq0$ para todo $(x,t)\in(\mathbb{R}^n\setminus\{0\})\times I$. Definimos entonces
\begin{align*}
H:(\mathbb{R}^n\setminus\{0\})\times I&\to\mathbb{R}^n\setminus\{0\}\\
(x,t)&\mapsto\frac{x}{\lVert x\rVert}t+(1-t)x.
\end{align*}
Como $H$ se forma a partir de funciones continuas mediante operaciones que preservan la continuidad, se tiene que $H$ es continua. Además
\begin{itemize}
\item Para cualquier $x\in\mathbb{R}^n\setminus\{0\}$, $H(x,0)=x$ y, como $\lVert H(x,1)\rVert=\left\lVert\frac{x}{\lVert x\rVert}\right\rVert=1$, entonces $H(x,1)\in S^{n-1}$.
\item Para cualesquiera $z\in S^{n-1}$ y $t\in I$, se tiene $H(z,t)=\frac{z}{\lVert z\rVert}t+(1-t)z=zt+(1-t)z=z$.
\end{itemize}
Por tanto $S^{n-1}$ es un retracto de deformación de $\mathbb{R}^n\setminus\{0\}$.
\end{proof}
\textbf{b)} $\mathbb{R}^3-\{(0,0,0)\}$.
\begin{proof}[\textbf{Solución}]
Por el \Cref{Lema:Esfera_retracto_de_Rn_menos_origen}, $S^2$ es un retracto de deformación de $\mathbb{R}^3\setminus\{0\}$. Por tanto, para cualquier $z_0\in S^2$ se tiene
\[
\pi_1(\mathbb{R}^3\setminus\{0\},z_0)\cong\pi_1(S^2,z_0)=\{*\}.
\]
Como $\mathbb{R}^3\setminus\{0\}$ es conexo por caminos, para cualquier $z\in\mathbb{R}^3\setminus\{0\}$ se tiene $\pi_1(\mathbb{R}^3\setminus\{0\},z)\cong\pi_1(\mathbb{R}^3\setminus\{0\},z_0)$ y por tanto $\pi_1(\mathbb{R}^3\setminus\{0\},z)\cong\{*\}$.
\end{proof}
